\Askhsh{35774}{2}
\begin{ylh}{1}
    \item 2.3 Μέτρα Θέσης και Διασποράς
\end{ylh}
Οι βαθμοί που έγραψε στα πέντε τεστ του 1\textsuperscript{ου} τετραμήνου, για το μάθημα των Μαθηματικών, ένας μαθητής της Γ' Λυκείου είναι: 12, 14, 8, 20, 16.
\begin{alist}
    \item Να βρείτε τη μέση τιμή $\bar{x}$ της βαθμολογίας του μαθητή, για τα πέντε τεστ. \monades{8}
    \item 
    \begin{rlist}
        \item Να συμπληρώσετε τα κενά στον παρακάτω πίνακα.
     \begin{center}
	\begin{mytblr}{}
	    Βαθμός $x_i$ & $x_i - \bar{x}$ & $(x_i - \bar{x})^2$ \\
	    12 & & \\
	    14 & & \\
	    8  & & \\
	    20 & & \\
	    16 & & \\
	    ΣΥΝΟΛΟ & \SetCell{bg=gray!25}{} & \\
	\end{mytblr}
\end{center}
        
        \monades{8}
        \item Αξιοποιώντας τον συμπληρωμένο πίνακα του ερωτήματος (β, i), να βρείτε τη διακύμανση και την τυπική απόκλιση, της βαθμολογίας του μαθητή για τα πέντε τεστ. \monades{9}
    \end{rlist}
\end{alist}

